\begin{abstract}
Dispatching rules such as earliest-deadline-first (EDF) and shortest-processing-time (SPT) assume constant processing capacity, which fails when the resource is a person whose attention depletes, whose fatigue accumulates and whose alertness follows a circadian rhythm. We formulate human-centric task scheduling as a partially observable Markov decision process in which rest is an explicit action, attention and fatigue are latent states observed through noisy estimates, and a multi-objective reward balances deadline-weighted task value against cognitive load. We propose an attention-aware deep Q-network (AA-DQN) that combines an observation history, as an approximation of the belief state, with Double Q-learning and a dueling architecture. In a cognitively grounded simulator, AA-DQN was benchmarked against nine rule-based schedulers, including grid-tuned EDF and SPT rules with periodic or state-triggered rest, using ten training seeds, 500 held-out workdays with common random numbers and multiplicity-corrected tests. Rules that never rest performed poorly (EDF: \edfOnTimePct\% of tasks on time). The pre-specified AA-DQN matched the return of the best tuned rule and completed more tasks on time (\aaOnTimePct\% versus \bestHeurOnTimePct\%). A compact variant with intensity-only actions, evaluated as a secondary analysis, exceeded the best rule in return (+\SecRetSPTThreshold, $p<0.001$) and on-time completion (\AADQNThreeactOnTimePct\%). The learned policies alone combined high throughput with protection of demanding tasks: they front-load hard work while attention is high and rest around the post-lunch dip. Their advantage grew markedly with accurate cognitive-state estimates, and they were robust to sensor noise, whereas generalisation to unseen workloads remained limited.
\end{abstract}

\begin{keyword}
Human-centric scheduling \sep Deep reinforcement learning \sep Partially observable Markov decision process \sep Cognitive fatigue \sep Attention \sep Rest-break scheduling \sep Industry 5.0
\end{keyword}
